\section{Project Feasibility}
\label{sec:feasibility}

\subsection{Technical Feasibility}
\begin{itemize}
    \item \textbf{IDEs:} Visual Studio Code, Android Studio
    \item \textbf{Tools \& Utilities:} GitHub, Postman, Figma, MySQL Workbench
    \item \textbf{Hardware:} Android smartphones, Development laptops/desktops, Cloud hosting server
\end{itemize}

\subsection{Technical Feasibility}
\subsubsection{IDEs}
\begin{itemize}
    \item Visual Studio Code
    \item Android Studio
\end{itemize}

\subsubsection{Tools \& Utilities}
\begin{itemize}
    \item GitHub
    \item Postman
    \item Figma
    \item pgAdmin / PostgreSQL Management Tools
\end{itemize}

\subsubsection{Hardware}
\begin{itemize}
    \item Android smartphones for field operations
    \item Development laptops/desktops
    \item Cloud hosting server
\end{itemize}

\subsubsection{Technologies}
\begin{itemize}
    \item Frontend: React Native
    \item Backend: Node.js with NestJS
    \item Database: PostgreSQL
    \item Authentication: JWT Authentication
    \item REST API Architecture
\end{itemize}

\subsection{Operational Feasibility}
The system is operationally feasible because it is designed according to the existing tea collection workflow currently practiced in Sri Lanka. The system does not attempt to completely replace the current operational process. Instead, it digitizes manual activities such as tea collection recording, weight verification, truck allocation, factory request handling, and monthly payment calculations. The user interface will be kept simple to support farmers, collection center employees, truck operators, and factory staff who may have limited technical knowledge. The system will reduce manual paperwork, improve communication efficiency, and minimize record-keeping errors.

\subsection{Legal \& Ethical Feasibility}
The proposed system is legally and ethically feasible because it focuses on secure and authorized management of operational data. The system will protect user information through secure login authentication, restrict unauthorized access using role-based access control, maintain secure records of transactions and payments, and ensure data confidentiality for farmers, factories, and employees. The project does not violate any known legal regulations and will operate within standard business and data management practices.

\subsection{Social Feasibility}
The system supports digital transformation within the Sri Lankan tea industry. The proposed system improves operational transparency, communication between stakeholders, record accuracy, and efficiency of tea distribution operations. The system can positively impact farmers, tea collection centers, factory operators, and transport coordinators.

\subsection{Economic Feasibility}
The proposed system is economically feasible because it reduces operational inefficiencies and minimizes manual administrative work. The system can help reduce paper-based record costs, manual calculation errors, data duplication, and time spent managing records. Since the project uses open-source technologies such as React Native, Node.js, NestJS, and PostgreSQL, software development costs are relatively low. In the long term, the system can improve operational productivity and reduce management expenses for tea distribution centers.

\subsection{Schedule Feasibility}
The project is feasible within the allocated academic project timeline. The initial Minimum Viable Product (MVP) can include user authentication, tea collection management, truck management, factory request handling, payment calculations, and reporting features. Additional advanced features can be added incrementally in future development phases. The selected technology stack supports rapid development, testing, and deployment, making the implementation achievable within the planned schedule.